\par\Needspace{6\baselineskip}\medskip
\begin{center}\hypertarget{native-186-article139}{}\pdfbookmark[2]{139.}{article-139}139.\end{center}
Es seien $C,C',C''$ die drei heliozentrischen Orte des Himmelskörpers auf der Himmelskugel. Sie
liegen der Reihe nach auf den Großkreisen $AB,A'B',A''B''$, und zwar zwischen $A$ und $B$, $A'$ und
$B'$, $A''$ und $B''$ (Artikel 64, III). Außerdem liegen die Punkte $C,C',C''$ auf demselben
Großkreis, nämlich auf dem Schnittkreis der Bahnebene mit der Himmelskugel. Wir
bezeichnen mit $r,r',r''$ die drei Abstände des Himmelskörpers von der Sonne, mit $\rho,\rho',\rho''$ seine
Abstände von der Erde und mit $R,R',R''$ die Abstände der Erde von der Sonne. Ferner setzen wir die
Bögen $C'C'',CC'',CC'$ der Reihe nach $=2f,2f',2f''$ und
\[
r'r''\sin2f=n,\qquad rr''\sin2f'=n',\qquad rr'\sin2f''=n''.
\]
\hypertarget{native-187}{}Damit gilt
\[
\begin{gathered}
f'=f+f''\\[.4ex]
AC+CB=\delta,\qquad A'C'+C'B'=\delta',\\[.4ex]
A''C''+C''B''=\delta'',
\end{gathered}
\]
und außerdem
\[
\begin{aligned}
\frac{\sin\delta}{r}&=\frac{\sin AC}{\rho}=\frac{\sin CB}{R}\\[.7ex]
\frac{\sin\delta'}{r'}&=\frac{\sin A'C'}{\rho'}=\frac{\sin C'B'}{R'}\\[.7ex]
\frac{\sin\delta''}{r''}&=\frac{\sin A''C''}{\rho''}=\frac{\sin C''B''}{R''}\mathpunct{\cdot}
\end{aligned}
\]

Daraus ist ersichtlich, dass sich, sobald die Lage der Punkte $C,C',C''$ bekannt ist, die Größen
$r,r',r'',\rho,\rho',\rho''$ bestimmen lassen. Wir werden nun zeigen, wie sich diese Lage aus den Größen
\[
\frac{n''}{n}=P,\qquad 2\left(\frac{n+n''}{n'}-1\right)r'^3=Q
\]
ableiten lässt, die wir oben als Ausgangspunkt unseres Verfahrens angegeben haben.
